% GENERATED FILE. Edit canonical lessons, then run npm run generate:books.
% canonical-source-hash: fnv1a64:517e9cf3e8bb6ca0
% canonical-lessons: TA-C255-cuvaci, TA-C255-kunamaku, TA-C255-vinku, TA-C255-eluppu, TA-C255-pal-tulakku

\chapter{To breathe, To recover, To swell, To wake, To brush one's teeth}
\label{ch:ta-to-breathe-to-recover-to-swell-to-wake-to-brush-one-s-teeth}
\hlchaptermodality{255}

\begin{chapteropening}
\textbf{By the end of this chapter:} \emph{I can say to breathe, to recover, to swell, to wake and to brush one's teeth.}

Complete the last lesson of chapter 255: I can say to breathe, to recover, to swell, to wake and to brush one's teeth.
\end{chapteropening}

\section[\texorpdfstring{\ta{சுவாசி} (cuvāci)}{சுவாசி (cuvāci)}]{\ta{சுவாசி} (cuvāci) — to breathe}

\label{lesson:TA-C255-cuvaci}



\begin{quote}
Before the new one: say the Tamil for to enter, then the Tamil for to exit.
\end{quote}

\subsection*{You'll want to know: \ta{சுவாசி}}
\textbf{\ta{சுவாசி}} — \emph{cuvāci} — \textquotedblleft{}to breathe\textquotedblright{}.

\begin{grammarlens}[title={What you've built}]
One more word for this chapter.
\end{grammarlens}

\subsection*{Your turn}
\begin{itemize}
\raggedright
  \item \emph{Say it:} \emph{cuvāci}
  \item \emph{Say it:} \emph{cuvāci}, once more
  \item \emph{Say it:} say \emph{cuvāci}
  \item \emph{Recall:} say the Tamil for to enter, then the Tamil for to exit, then say \emph{cuvāci} again
\end{itemize}

\begin{tcolorbox}[breakable,colback=teal!4,colframe=teal!35!black,title={Before you move on}]
What does \ta{சுவாசி} mean? (\textquotedblleft{}To breathe\textquotedblright{}.) Say it once more.
\end{tcolorbox}
\section[\texorpdfstring{\ta{குணமாகு} (kuṇamāku)}{குணமாகு (kuṇamāku)}]{\ta{குணமாகு} (kuṇamāku) — to recover, to get well}

\label{lesson:TA-C255-kunamaku}



\begin{quote}
Before the new one: say the Tamil for to exit, then the Tamil for to breathe.
\end{quote}

\subsection*{You'll want to know: \ta{குணமாகு}}
\textbf{\ta{குணமாகு}} — \emph{kuṇamāku} — \textquotedblleft{}to recover, to get well\textquotedblright{}.

\begin{grammarlens}[title={What you've built}]
One more word for this chapter.
\end{grammarlens}

\subsection*{Your turn}
\begin{itemize}
\raggedright
  \item \emph{Say it:} \emph{kuṇamāku}
  \item \emph{Say it:} \emph{kuṇamāku}, once more
  \item \emph{Say it:} say \emph{kuṇamāku}
  \item \emph{Recall:} say the Tamil for to exit, then the Tamil for to breathe, then say \emph{kuṇamāku} again
\end{itemize}

\begin{tcolorbox}[breakable,colback=teal!4,colframe=teal!35!black,title={Before you move on}]
What does \ta{குணமாகு} mean? (\textquotedblleft{}To recover, to get well\textquotedblright{}.) Say it once more.
\end{tcolorbox}
\section[\texorpdfstring{\ta{வீங்கு} (vīṅku)}{வீங்கு (vīṅku)}]{\ta{வீங்கு} (vīṅku) — to swell}

\label{lesson:TA-C255-vinku}



\begin{quote}
Before the new one: say the Tamil for to breathe, then the Tamil for to recover.
\end{quote}

\subsection*{You'll want to know: \ta{வீங்கு}}
\textbf{\ta{வீங்கு}} — \emph{vīṅku} — \textquotedblleft{}to swell\textquotedblright{}.

\begin{grammarlens}[title={What you've built}]
One more word for this chapter.
\end{grammarlens}

\subsection*{Your turn}
\begin{itemize}
\raggedright
  \item \emph{Say it:} \emph{vīṅku}
  \item \emph{Say it:} \emph{vīṅku}, once more
  \item \emph{Say it:} say \emph{vīṅku}
  \item \emph{Recall:} say the Tamil for to breathe, then the Tamil for to recover, then say \emph{vīṅku} again
\end{itemize}

\begin{tcolorbox}[breakable,colback=teal!4,colframe=teal!35!black,title={Before you move on}]
What does \ta{வீங்கு} mean? (\textquotedblleft{}To swell\textquotedblright{}.) Say it once more.
\end{tcolorbox}
\section[\texorpdfstring{\ta{எழுப்பு} (eḻuppu)}{எழுப்பு (eḻuppu)}]{\ta{எழுப்பு} (eḻuppu) — to wake (someone) up}

\label{lesson:TA-C255-eluppu}



\begin{quote}
Before the new one: say the Tamil for to recover, then the Tamil for to swell.
\end{quote}

\subsection*{You'll want to know: \ta{எழுப்பு}}
\textbf{\ta{எழுப்பு}} — \emph{eḻuppu} — \textquotedblleft{}to wake (someone) up\textquotedblright{}.

\begin{grammarlens}[title={What you've built}]
One more word for this chapter.
\end{grammarlens}

\subsection*{Your turn}
\begin{itemize}
\raggedright
  \item \emph{Say it:} \emph{eḻuppu}
  \item \emph{Say it:} \emph{eḻuppu}, once more
  \item \emph{Say it:} say \emph{eḻuppu}
  \item \emph{Recall:} say the Tamil for to recover, then the Tamil for to swell, then say \emph{eḻuppu} again
\end{itemize}

\begin{tcolorbox}[breakable,colback=teal!4,colframe=teal!35!black,title={Before you move on}]
What does \ta{எழுப்பு} mean? (\textquotedblleft{}To wake (someone) up\textquotedblright{}.) Say it once more.
\end{tcolorbox}
\section[\texorpdfstring{\ta{பல்} \ta{துலக்கு} (pal tulakku)}{பல் துலக்கு (pal tulakku)}]{\ta{பல்} \ta{துலக்கு} (pal tulakku) — to brush one's teeth}

\label{lesson:TA-C255-pal-tulakku}



\begin{quote}
Before the new one: say the Tamil for to swell, then the Tamil for to wake.
\end{quote}

\subsection*{You'll want to know: \ta{பல்} \ta{துலக்கு}}
\textbf{\ta{பல்} \ta{துலக்கு}} — \emph{pal tulakku} — \textquotedblleft{}to brush one's teeth\textquotedblright{}.

\begin{grammarlens}[title={What you've built}]
One more word for this chapter.
\end{grammarlens}

\subsection*{Your turn}
\begin{itemize}
\raggedright
  \item \emph{Say it:} \emph{pal tulakku}
  \item \emph{Say it:} \emph{pal tulakku}, once more
  \item \emph{Say it:} say \emph{pal tulakku}
  \item \emph{Recall:} say the Tamil for to swell, then the Tamil for to wake, then say \emph{pal tulakku} again
\end{itemize}

\begin{tcolorbox}[breakable,colback=teal!4,colframe=teal!35!black,title={Before you move on}]
What does \ta{பல்} \ta{துலக்கு} mean? (\textquotedblleft{}To brush one's teeth\textquotedblright{}.) Say it once more.
\end{tcolorbox}
